\section{Equações do Movimento Orbital e Manobras Associadas ao Problema de RVD/B}
\label{sec:eq_mov_orbital_manobras}

A formulação da dinâmica orbital no contexto de missões de RVD/B leva em consideração as diferentes fases da missão.
Quando o lançamento atinge a altitude nominal, é realizada a injeção do perseguidor em órbita.
Esses eventos levam em consideração os algoritmos previstos para a injeção em uma órbita adequada para iniciar as manobras relacionadas ao \gls{rvdb}.
Neste trabalho, assume-se que o perseguidor  foi injetado em uma órbita estável, suficientemente próxima do alvo, de forma que a órbita de fase se inicie.
O sistema de controle do lançador possui as entradas (dados) necessários para assegurar a injeção do perseguidor em uma órbita coplanar à do alvo, e ao ajuste do ângulo de fase para que ocorra a transição às etapas de encontro.

\subsection{Lançamento: vínculos associados às missões de RVD/B}
\label{sec:lancamento_vinculos}

O início da missão de RVD/B ocorre na etapa de lançamento, conforme ilustrado na Figura~\ref{fig:etapas_rvdb}.
Nessa fase, a espaçonave perseguidora deve ser injetada em uma órbita compatível com a do alvo, usando a janela de lançamento planejada e evitar inserções em órbitas inadequadas, que poderão demandar correções adicionais de trajetória.  

\begin{figure}[ht]
\centering
\includegraphics[width=1\linewidth]{3 Formulacao matematica/Imagens/Lançamento quadro explicativo.png}
\caption{Fases de uma missão de RVD/B, adaptado de \cite{fehse2003}.}
\label{fig:etapas_rvdb}
\end{figure}

A definição da janela de lançamento é vinculada à órbita conhecida do alvo. Os parâmetros de injeção são escolhidos em função dessa órbita, de forma a garantir a inserção do veículo em uma trajetória coplanar à do alvo.

\subsubsection{Detalhes sobre a geometria da órbita do perseguidor}

Para que o veículo lançador insira a espaçonave perseguidora na órbita de injeção desejada, é necessário fornecer os elementos orbitais \cite{vallado2007fundamentals} de injeção ou o vetor posição--velocidade no referencial inercial (ECI) .
Neste trabalho foram considerados os elementos orbitais clássicos $(a, e, i, \Omega, \omega, M)$ e o vetor posição no referencial inercial no instante de injeção orbital.
Além disso, as condições para que a injeção orbital seja considerada válida incluem:

\begin{enumerate}[label=(\roman*)]
    \item o correto posicionamento no semi-eixo maior da órbita inicial do perseguidor;
    \item fase orbital correta no momento da separação;
\end{enumerate}

% \cite{bate1971fundamentals, vallado2007fundamentals, curtis2010orbital}.
